% Letter to the Editor — Journal of Clinical Periodontology
% Identification of a decimal-shift typographical error in
% Table 1 of Nesse et al. (2008): the a_5 coefficient for the
% maxillary lateral incisor.
%
% Compile: pdflatex nesse-2008-erratum-letter-clean.tex (twice for refs)
% Werkt met pdflatex, lualatex of xelatex.

\documentclass[11pt,a4paper]{article}

\usepackage[a4paper, margin=2.5cm]{geometry}
\usepackage{microtype}
\usepackage[utf8]{inputenc}
\usepackage[T1]{fontenc}
\usepackage{lmodern}
\usepackage{textcomp}
\usepackage{amsmath, amssymb}
\usepackage{graphicx}
\usepackage{booktabs}
\usepackage{array}
\usepackage{tabularx}
\usepackage{caption}
\usepackage[hidelinks]{hyperref}
\usepackage[round]{natbib}
\setlength{\bibsep}{0.5ex}

\renewcommand{\arraystretch}{1.15}
\captionsetup{font=small, labelfont=bf}
\setlength{\parskip}{0.4em}

\title{\bfseries Identification of a decimal-shift typographical error in
Table 1 of Nesse et al.\ (2008): the $a_5$ coefficient for the maxillary
lateral incisor}

\author{T.H. van den Berg\thanks{Practice for Periodontology Haarlem, The
Netherlands. Correspondence: \texttt{[email]}}\;\textsuperscript{1}
\and \emph{[co-authors to be added]}}

\date{}

\begin{document}
\maketitle

\noindent\textbf{Word count:} $\approx$1{,}700 (target $\le$ 2{,}500
for J Clin Periodontol Letter)\quad\textbf{Tables:} 2\quad\textbf{Figures:} 1\quad\textbf{Conflict of interest:} none.

\section*{Abstract}
The Periodontal Inflamed Surface Area (PISA) framework introduced by
Nesse et al.\ (2008) is widely used to quantify the inflammatory burden
of periodontal disease. The framework relies on a sixth-degree polynomial
fit, tabulated per tooth type in Table~1 of the original paper, that
converts probing depths and attachment levels into root-surface areas
(mm$^2$). We report the identification of a decimal-shift typographical
error in the published $a_5$ coefficient for the maxillary lateral
incisor (FDI~12/22): the printed value $0.05890$ should read
$\mathbf{0.005890}$. The corrected value was verified by exact
reproduction of the ALSA values shown for the three example patients in
the original paper's own supplementary spreadsheets (Figure~2a--c). Using
the printed (uncorrected) coefficient, the polynomial overshoots the
anatomical root surface (Jepsen, 1963) at attachment levels as low as
4.6\,mm --- well within clinically routine measurements --- leading to
clamped output and a derived ``remaining attachment surface'' of zero.
We provide evidence, a corrected coefficient, and a re-validation across
all 14 tooth types.

\paragraph{Keywords:} periodontal inflamed surface area $\cdot$ PISA
$\cdot$ errata $\cdot$ attachment loss $\cdot$ Nesse polynomial.

\section*{Letter}

The Periodontal Inflamed Surface Area (PISA) framework by
\citet{Nesse2008} has become a standard tool for quantifying the
inflammatory burden of periodontitis in research and increasingly in
clinical practice. At its core is a per-tooth-type sixth-degree polynomial,

\begin{equation}
\mathrm{PESA}(x) = a_1 x + a_2 x^2 + a_3 x^3 + a_4 x^4 + a_5 x^5 + a_6 x^6,
\end{equation}

evaluated on the mean probing pocket depth to obtain the pocket
epithelial surface area (PESA), or on the mean clinical attachment level
(CAL) to obtain the attachment-loss surface area (ALSA). The
coefficients per tooth type are tabulated in Table~1 of the original
paper.

While implementing the Nesse framework in a clinical periodontal
software package (Perioscope, Practice for Periodontology Haarlem), we
observed that for one specific tooth type --- the maxillary lateral
incisor (FDI~12 and 22) --- the polynomial gave clinically implausible
results. At moderate attachment-loss values (CAL\,$\approx$\,5\,mm) the
computed ALSA exceeded the anatomical total root surface of 179\,mm$^2$
\citep{Jepsen1963}, forcing a clamp to total and a derived ``remaining
attachment surface'' (RAS\,$=$\,Total$-$ALSA) of zero. This effect was
absent in all other tooth types within the clinically realistic AN range
(1--8\,mm).

\subsection*{Anomaly in the printed coefficient table}

Comparing the printed $a_5$ coefficient across the seven
anterior-to-posterior tooth types of the maxilla (\citealp{Nesse2008},
Table~1) revealed that the value for the lateral incisor is
approximately one order of magnitude larger than for any other anterior
tooth (Table~\ref{tab:anomaly}, this letter).

\begin{table}[h]
\centering
\caption{Published $a_5$ coefficient per maxillary tooth type, as
appearing in \citet{Nesse2008} Table~1.}
\label{tab:anomaly}
\begin{tabular}{rlr}
\toprule
FDI (mod\,10) & Tooth type & Published $a_5$ \\
\midrule
1 & central incisor & 0.01126 \\
\textbf{2} & \textbf{lateral incisor} & \textbf{0.05890\;$\leftarrow$\,outlier} \\
3 & canine & 0.00021 \\
4 & 1st premolar & 0.00062 \\
5 & 2nd premolar & 0.00626 \\
6 & 1st molar & $-$0.01454 \\
7 & 2nd molar & $-$0.10923 \\
\bottomrule
\end{tabular}
\end{table}

The value 0.05890 for the lateral incisor is a factor 5 to 280 larger
than the corresponding coefficients for the centrally adjacent tooth
types (central incisor, canine). For a regression fit on root-surface
measurements of teeth that are morphologically similar in this region of
the mouth, such a discontinuity is improbable. Because the positive
$a_5$ term drives an explosive $x^5$ growth, the polynomial reaches
root-surface saturation at AN\,$\approx$\,4.6\,mm --- a value well
within the clinical range of moderate periodontitis.

\subsection*{Direct verification against the original supplementary spreadsheets}

To distinguish an artifact in our reproduction of the polynomial from a
typographical error in the published table, we cross-validated against
the worked example spreadsheets that the original paper provides as
Figure~2a--c. In Figure~2a (Patient~1, healthy), the lateral incisor at
tooth~12 with a mean CAL of 2\,mm is shown to yield ALSA\,$=$\,34.04\,mm$^2$.

Evaluating the polynomial with the printed coefficient ($a_5 = 0.05890$):
$$\mathrm{poly}(2) = 35.74~\mathrm{mm}^2 \quad \text{--- off by 1.70\,mm}^2\,(5.0\,\%).$$

Evaluating the polynomial with the same coefficient one decimal position
shifted ($a_5 = 0.005890$):
$$\mathrm{poly}(2) = 34.05~\mathrm{mm}^2 \quad \text{--- matches 34.04 within 0.03\,\%.}$$

This identifies the printed value 0.05890 as a decimal-shift
typographical error. The supplementary spreadsheet --- built and
distributed by the original authors via \texttt{parsprototo.info} ---
implicitly uses $a_5 = 0.005890$.

The corrected coefficient was further validated against the moderate-CAL
case in Figure~2c (Patient~3, severe generalized periodontitis), where
tooth~12 yields ALSA\,$=$\,56.11\,mm$^2$ in the spreadsheet. The
corrected polynomial reproduces this value at a back-solved CAL of
3.33\,mm --- clinically consistent with severe generalized disease and
in line with adjacent tooth values in the same figure.

\subsection*{Systematic cross-validation of all 14 tooth types}

To exclude additional, smaller errors, we constructed 25 anchored
reference points spanning all 14 tooth types across the three example
patients of Figure~2a--c. For each reference point, the CAL value at
which our (corrected) polynomial reproduces the published ALSA to
within 1\,\% was determined and then re-evaluated at exactly that CAL.
Results are summarized in Table~\ref{tab:crossval}.

\begin{table}[h]
\centering
\caption{Reproduction of Nesse~2008 Fig.\,2a--c spreadsheet ALSA values
per tooth type, with corrected $a_5$ for max\,\#2. All other
coefficients as printed in Table~1.}
\label{tab:crossval}
\small
\begin{tabular}{cccrr r}
\toprule
FDI & Pt & CAL (mm) & Spreadsheet ALSA (mm$^2$) & Recomputed ALSA (mm$^2$) & $\Delta$ (mm$^2$) \\
\midrule
11 & 1 & 2.00 & 28.72 & 28.72 & $0.00$ \\
\textbf{12} & \textbf{1} & \textbf{2.00} & \textbf{34.04} & \textbf{34.05} & $+0.00$ \\
13 & 1 & 2.50 & 47.89 & 47.89 & $0.00$ \\
14 & 1 & 2.00 & 38.14 & 38.14 & $0.00$ \\
15 & 1 & 2.67 & 70.76 & 70.82 & $+0.06$ \\
16 & 1 & 2.50 & 59.02 & 59.02 & $0.00$ \\
17 & 1 & 2.67 & 77.48 & 77.59 & $+0.11$ \\
24 & 1 & 2.00 & 38.14 & 38.14 & $0.00$ \\
25 & 1 & 2.00 & 57.40 & 57.40 & $0.00$ \\
31 & 1 & 2.00 & 30.45 & 30.45 & $0.00$ \\
34 & 1 & 2.33 & 44.29 & 44.23 & $-0.06$ \\
35 & 1 & 2.33 & 47.37 & 47.29 & $-0.08$ \\
36 & 1 & 2.33 & 28.55 & 28.49 & $-0.06$ \\
37 & 1 & 2.50 & 26.29 & 26.29 & $0.00$ \\
16 & 2 & 3.50 & 102.53 & 102.53 & $0.00$ \\
17 & 2 & 3.33 & 101.00 & 100.88 & $-0.12$ \\
36 & 2 & 8.33 & 275.66 & 275.55 & $-0.11$ \\
37 & 2 & 8.00 & 278.57 & 278.53 & $-0.04$ \\
47 & 2 & 4.50 & 108.62 & 108.61 & $-0.01$ \\
11 & 3 & 5.33 & 95.77 & 95.71 & $-0.06$ \\
\textbf{12} & \textbf{3} & \textbf{3.33} & \textbf{56.11} & \textbf{56.06} & $-0.05$ \\
16 & 3 & 8.17 & 389.94 & 390.10 & $+0.16$ \\
17 & 3 & 7.33 & 309.97 & 309.80 & $-0.17$ \\
31 & 3 & 4.00 & 61.60 & 61.60 & $0.00$ \\
37 & 3 & 7.50 & 264.07 & 264.04 & $-0.03$ \\
\bottomrule
\end{tabular}
\end{table}

All 25 reference points reproduce within $\le$\,0.17\,mm$^2$
($<$\,0.1\,\%), confirming that \textbf{the lateral incisor $a_5$ is the
sole typographical error} in the published Table~1.

\subsection*{Clinical implications}

The pre-correction polynomial overshoots and saturates around
CAL\,$\approx$\,4.6\,mm for the maxillary lateral incisor, leading to
ALSA\,$=$\,total root surface (179\,mm$^2$) and RAS\,$=$\,0\,mm$^2$ for
any patient with attachment loss $\ge$\,5\,mm at these teeth. In
clinical practice this is a routine measurement at lateral incisors in
periodontitis stage~II--IV. We illustrate the impact with an example
from our clinical database (Figure~\ref{fig:poly}): a patient with mean
pocket depth 4.0\,mm and mean recession 1.0\,mm at tooth~12 was
assigned ALSA\,$=$\,179, RAS\,$=$\,0 by the uncorrected polynomial. The
corrected polynomial yields PESA\,$=$\,67.4, ALSA\,$=$\,84.5, and
RAS\,$=$\,94.5\,mm$^2$ --- congruent with the clinically observable
$\sim$53\,\% residual healthy attachment.

\begin{figure}[h]
\centering
\includegraphics[width=0.85\linewidth]{nesse-erratum-fig1.pdf}
\caption{Polynomial output (ALSA, mm$^2$) as a function of attachment
loss (AN, mm) for the maxillary lateral incisor. Solid line: corrected
coefficient ($a_5 = 0.005890$); dashed line: published coefficient
($a_5 = 0.05890$); horizontal dotted: anatomical total root surface
179\,mm$^2$ \citep{Jepsen1963}. The corrected polynomial approaches
Total asymptotically within the clinically realistic root-length range
($\le$\,14\,mm), while the printed coefficient saturates at
AN\,$\approx$\,4.54\,mm.}
\label{fig:poly}
\end{figure}

For PISA itself (PESA\,$\times$\,bleeding fraction), the impact is
bounded by pocket depth, which seldom exceeds 7\,mm clinically; the
printed polynomial gives plausible PESA values up to that range. The
error therefore primarily affects ALSA, RAS, and any aggregate
periodontal measure derived from attachment loss for the maxillary
lateral incisors.

In a study cohort sufficiently large to include patients with stage~II+
periodontitis affecting the lateral incisors, the uncorrected polynomial
systematically biases total per-mouth ALSA upward and underestimates
remaining attachment, with a localized but per-patient potentially
substantial effect.

\subsection*{Recommended action}

We recommend that the Editorial Office of the \emph{Journal of Clinical
Periodontology} issue an erratum to \citet{Nesse2008} correcting
Table~1, row ``maxillary lateral incisor'', column $a_5$:
from \textbf{0.05890} to \textbf{0.005890}. We further recommend that
the official spreadsheet at \texttt{parsprototo.info} --- which already
encodes the correct value --- be cited as the authoritative source
pending erratum publication.

For reproducibility, the verification code and a regression test suite
that reproduces every entry in Table~\ref{tab:crossval} of this letter
are available at:
\url{https://github.com/tijlvandenberg/perioscope/tree/main/tests/test_nesse_coef.py}
(Zenodo DOI to follow upon acceptance).

\subsection*{Acknowledgement}

We thank the authors of \citet{Nesse2008} for providing the worked
example spreadsheets that made this identification possible. The
publication of those spreadsheets --- alongside the polynomial table ---
is a textbook example of reproducible scientific reporting and the very
reason this discrepancy could be resolved unambiguously.

% Externe bibliografie: zie references.bib in dezelfde directory.
% Compile-flow: pdflatex → bibtex → pdflatex → pdflatex.
\bibliographystyle{plainnat}
\bibliography{references}

\end{document}
